\chapter{第八章：Agent 架构模式}

很多工程师第一次做 Agent，实际上只是把一个大语言模型（Large Language Model, LLM）包装成``问一句、答一句''的聊天接口。这样的系统当然有价值；但一旦要让模型读文件、调用 API、执行命令或修改外部状态，问题就不再是``模型够不够聪明''，而是\textbf{谁拥有执行权、谁保存状态、谁在失控前踩刹车}。

因此不要把 Agent 想成``更自主的聊天机器人''。更可操作的定义是：

\begin{quote}
\textbf{Agent 是一个宿主 runtime：模型只在其中提出下一步意图，runtime 负责验证、执行、记录结果并决定何时停止。}
\end{quote}

规划、长期记忆、多 Agent 都是可选能力，不是 Agent 的入场券。一个只有``模型选择工具 → 宿主执行 → 结果回填 → 下一轮''的受限循环，已经是 Agent；一个带 Planner 但没有权限边界和停止条件的系统，仍然不可靠。


\section{8.1 什么样的系统才算 Agent}

先看最小对比：


\begingroup
\scriptsize
\setlength{\tabcolsep}{3pt}
\begin{longtable}{|p{0.21\textwidth}|p{0.22\textwidth}|p{0.24\textwidth}|p{0.23\textwidth}|}
\hline
运行时形态 & 模型能决定什么 & 宿主必须保证什么 & 典型边界 \\ \hline
单次 LLM 调用 & 只生成文本 & 输入输出格式、超时、内容安全 & 没有外部副作用 \\ \hline
一次工具调用 & 选择工具和参数 & schema 校验、权限、工具结果回填 & 不能把 tool call 当已执行 \\ \hline
受限 Agent loop & 每轮决定“回答或继续调用” & 最大步数、token 预算、sandbox、审计 & 最适合从零落地 \\ \hline
显式工作流 / graph & 只在允许节点做局部决策 & 状态机、审批、幂等、补偿 & 高风险或流程稳定的任务 \\ \hline
Planner + executor & 先生成计划，再决定局部重规划 & checkpoint、计划版本、执行预算 & 长任务，不是默认配置 \\ \hline
\end{longtable}
\endgroup

以本地 \texttt{agent-cli} 为例，\texttt{src/agent.ts} 的 \texttt{runAgent()} 并不让模型直接拥有文件或 shell 权限。它接收 \texttt{ModelProvider}、\texttt{Tool[]}、消息历史和 sandbox 配置；模型返回文本或 \texttt{toolCalls}，runtime 再决定是否执行。默认循环最多 10 轮。这个结构比``自主性强/弱''的说法更重要，因为它直接决定能否测试、审计和收紧权限。

做一个最小 Agent 前，先回答四个具体问题，而不是先背四个抽象名词：
\begin{enumerate}[leftmargin=2em]
\item \textbf{状态是什么？} 当前消息、工具结果、文件摘要和计划分别放在哪里，哪些能进入下一轮。
\item \textbf{模型能提什么意图？} 只允许文本，还是允许有限集合中的结构化 tool call。
\item \textbf{谁做最终校验？} 参数、目录、命令、URL、写操作审批不能交给模型自觉。
\item \textbf{如何停止？} 最大轮数、token 预算、用户取消、工具超时和不可恢复错误必须由 runtime 处理。
\end{enumerate}

如果你把``模型选择下一步''理解成概率决策，把``工具调用''理解成受策略约束的 API 编排，把``记忆''理解成状态管理，那么 Agent 架构其实就是\textbf{传统软件工程 + 概率模型决策层}。这也是为什么后端、平台、基础设施工程师转做 Agent 很有优势。


\section{8.2 Agent Loop：Perceive -> Reason -> Act -> Observe}

绝大多数 Agent 架构，本质上都是一个循环（loop），但不要把 \texttt{Reason} 误解为必须展示或保存完整思维链。工程上可观测、可控制的是：模型给出了文本还是结构化动作意图，runtime 是否批准，以及环境返回了什么。


\begin{codeblock}
messages + workspace context
            │
            ▼
    ModelProvider.complete(messages, tools)
            │
            ├── 无 tool call ──> final answer
            ▼
      tool call intent
            │
            ▼
 sandbox / policy validation ──拒绝──> tool error 回填
            │ 批准
            ▼
 Tool.execute(arguments)
            │
            ▼
 role: tool result 写回 messages ──> 下一轮
\end{codeblock}

在 \texttt{agent-cli} 中，这条路径分别落在 \texttt{cli.ts → provider.ts → agent.ts → sandbox.ts → tools/* → agent.ts}。这比 Perceive / Reason / Act / Observe 的术语更值得背，因为每一跳都能放测试、日志和权限检查。

\begin{itemize}[leftmargin=2em]
\item \textbf{Perceive}：送进模型的 messages、工作区摘要、检索结果；常见失败是上下文缺失、陈旧、越权或超预算。
\item \textbf{Reason}：文本输出或 \texttt{toolCalls}；常见失败是工具选错、参数不完整、无意义循环。
\item \textbf{Act}：经过 policy 后的 \texttt{Tool.execute()}；常见失败是权限越界、超时、非幂等重试。
\item \textbf{Observe}：规范化的 tool result 回填；常见失败是原始报错过长、结果丢失、无法继续决策。
\end{itemize}

关键边界是：\textbf{模型负责提出意图，runtime 负责执行权、状态写入和终止权。} Workflow 的路径多为预先编码；Agent 的工具选择可动态生成；Graph 则把两者混合，让模型只能在受控节点做局部选择。


\section{8.3 核心架构模式一：ReAct（Reason + Act）}

ReAct 是最经典、最容易落地的 Agent 模式。它不要求把 Thought 暴露给用户，也不要求模型输出固定的 ``Thought / Action'' 文本。其工程内核只是：\textbf{模型依据上一次观察提出下一步工具意图，宿主执行后把观察结果送回去。}


\subsection{8.3.1 用真实 tool protocol 看 ReAct}

在 \texttt{agent-cli} 里，一次``定位 TypeScript 错误''的 trace 更接近下面这样：


\begin{codeblock}
用户：运行类型检查并说明最先要修什么
模型：tool_call { name: "run_typecheck", arguments: {} }
runtime：检查 sandbox；执行 npm run typecheck
工具：返回 { success, errors: [{ file, line, column, code, message }] }
runtime：把结构化结果作为 role: tool 消息回填
模型：要么调用 read_file / semantic_search 获取上下文，要么输出修复建议
\end{codeblock}

这里没有必要暴露模型内部推理。你要记录的是 tool name、经过校验的参数、耗时、结果摘要、是否重试和最终结论。这样既能调试，也不会把不可靠的自由文本思考当成系统事实。


\begin{enumerate}[leftmargin=2em]
\item \textbf{可解释}：可以回放每步 Thought / Action / Observation。
\item \textbf{易调试}：工具失败时能看到是模型选错工具、参数不对，还是外部系统有问题。
\item \textbf{泛化强}：只要工具描述合理，模型能把不同任务映射到统一执行循环。
\end{enumerate}

\subsection{8.3.2 一个完整示例}

用户目标：``帮我判断一篇 AI 新闻是否值得转发到团队群，并给出一句摘要。''


ReAct trace 可以设计成：


\begin{codeblock}
Thought: 我需要先获取新闻内容。
Action: fetch_url
Action Input: {"url": "https://example.com/ai-news"}
Observation: 文章讲的是某模型发布新推理 API，包含价格与延迟数据。

Thought: 我需要提取关键信息：发布时间、能力提升、成本变化、是否影响工程实践。
Action: extract_facts
Action Input: {"text": "..."}
Observation: 提取到 4 个核心事实。

Thought: 用户关心``是否值得转发''，我需要从新颖性、可信度、与团队相关性三方面评估。
Action: none
Observation: none

Final Answer: 值得转发。原因是……一句摘要：……
\end{codeblock}

注意，ReAct 中的 Thought 不一定需要向最终用户暴露。生产环境通常会：


\begin{itemize}[leftmargin=2em]
\item 对用户隐藏完整推理；
\item 但在日志中保留结构化 trace；
\item 只展示必要的 tool call 和 final answer。
\end{itemize}

\subsection{8.3.3 ReAct 的局限}

ReAct 并不万能。典型问题包括：


\begin{itemize}[leftmargin=2em]
\item 复杂任务时，模型容易一步一步试错，成本高；
\item 循环次数过多时，状态膨胀，token 成本上升；
\item 如果工具很多，模型会``犹豫''或乱选工具；
\item 缺少显式长程规划，做项目级任务时效率偏低。
\end{itemize}

因此实际系统里会把 ReAct 与规划器组合，形成 Plan-and-Execute 或 LangGraph 式状态图；但不要为了``像 Agent''而强行加 Planner。


\section{8.4 核心架构模式二：Function Calling}

Function Calling（函数调用）不是完整 Agent 架构，但它是现代 Agent 最重要的基础协议之一。它允许模型不再输出``请你调用某个函数''，而是直接生成结构化调用请求。


\subsection{8.4.1 OpenAI / Anthropic 的基本思想}

核心思路相同：开发者把工具列表和参数模式（schema）发给模型，模型在生成时可选择：


\begin{enumerate}[leftmargin=2em]
\item 直接返回文本；
\item 返回一个或多个工具调用；
\item 等工具结果回来后再继续回答。
\end{enumerate}

典型工具定义包含三部分：


\begingroup
\scriptsize
\setlength{\tabcolsep}{3pt}
\begin{longtable}{|p{0.31\textwidth}|p{0.31\textwidth}|p{0.31\textwidth}|}
\hline
字段 & 作用 & 示例 \\ \hline
name & 工具唯一标识 & \texttt{get\_weather} \\ \hline
description & 告诉模型何时使用 & ``获取某城市未来 3 天天气预报'' \\ \hline
parameters schema & 参数约束 & JSON Schema / Pydantic model \\ \hline
\end{longtable}
\endgroup

\subsection{8.4.2 JSON Schema 示例}

\begin{codeblock}
{
  "name": "get_weather",
  "description": "获取指定城市未来最多7天的天气预报，适用于用户询问天气、温度、降雨概率",
  "parameters": {
    "type": "object",
    "properties": {
      "city": {
        "type": "string",
        "description": "城市名称，例如 Beijing 或 Shanghai"
      },
      "days": {
        "type": "integer",
        "minimum": 1,
        "maximum": 7,
        "description": "查询天数"
      }
    },
    "required": ["city", "days"],
    "additionalProperties": false
  }
}
\end{codeblock}

模型看到这个 schema 后，会尽量输出类似：


\begin{codeblock}
{
  "tool_name": "get_weather",
  "arguments": {
    "city": "Shanghai",
    "days": 3
  }
}
\end{codeblock}

\subsection{8.4.3 Function Calling 的价值}

\begin{itemize}[leftmargin=2em]
\item 大幅降低解析自然语言 Action 的脆弱性；
\item 参数验证可以交给 schema 层；
\item 更利于并行工具调用；
\item 更适合接企业 API、数据库、内部服务。
\end{itemize}

当面试官问``Function Calling 和 Prompt-based tool use 的区别是什么''，你可以回答：\textbf{前者把工具调用从文本协议升级成结构化协议，使推理和执行之间的接口更稳定、可验证、可监控。}


\section{8.5 核心架构模式三：Plan-and-Execute}

Plan-and-Execute（规划后执行）适合中长任务。它把系统拆成两个角色。需要强调的是：\textbf{它不是 ReAct 的必选升级。} 当前 \texttt{agent-cli} 的 runtime 有受限 loop、预算和上下文压缩能力，但还没有把 Planner 接入默认 CLI；这是合理的阶段性边界，而不是``少了一个 Agent 能力''。


\begin{enumerate}[leftmargin=2em]
\item \textbf{Planner（规划器）}：先给出步骤列表；
\item \textbf{Executor（执行器）}：逐步执行每一步，必要时再局部调整。
\end{enumerate}

ASCII 示意图：


\begin{codeblock}
User Goal
   |
   v
+-----------+
|  Planner  | ---> Step 1, Step 2, Step 3...
+-----------+
   |
   v
+-----------+    tool calls / retries / state updates
| Executor  | ---------------------------------------->
+-----------+
   |
   v
 Final Result
\end{codeblock}

\subsection{8.5.1 为什么复杂任务更适合这样做}

假设任务是：``分析竞争对手最近 30 天发布的 AI 产品动态，给出对我们路线图的建议。''


如果纯 ReAct，模型可能会边搜边想，十几步后还没有全局结构；而 Plan-and-Execute 可以先规划：


\begin{enumerate}[leftmargin=2em]
\item 找出竞争对手名单；
\item 抓取最近 30 天新闻与产品更新；
\item 提取发布内容中的能力、价格、目标用户；
\item 横向比较；
\item 结合我方路线图输出建议。
\end{enumerate}

优点：


\begin{itemize}[leftmargin=2em]
\item 更适合长任务；
\item 执行器上下文更聚焦；
\item 可以对计划做审计、缓存与人工审批；
\item 易于插入 checkpoint。
\end{itemize}

缺点：


\begin{itemize}[leftmargin=2em]
\item 计划质量决定上限；
\item 计划可能过时，需要中途重规划（re-plan）；
\item 额外增加一次 LLM 调用。
\end{itemize}

\section{8.6 核心架构模式四：Reflexion}

Reflexion（反思）可以理解为``把可验证失败转化为下一次输入''。它不是单独替代 ReAct，也不是让模型写一段``我以后会更小心''的空泛自评。当前 \texttt{agent-cli} 已有 TypeScript 错误解析和错误摘要，正好是 Reflexion 可以消费的事实基础；自动把错误检索、修复、验证和重试串成闭环仍是待完成工作。


典型流程：


\begin{codeblock}
执行任务 -> 检查结果 -> 判断是否满足目标 -> 若不满足则生成改进策略 -> 再执行
\end{codeblock}

例如写 SQL：


\begin{enumerate}[leftmargin=2em]
\item 模型先生成 SQL；
\item 执行报错：\texttt{column usernmae does not exist}；
\item Reflexion 模块总结：``我拼错了列名 username''；
\item 下一轮带着纠错经验重试。
\end{enumerate}

Reflexion 的收益在于：


\begin{itemize}[leftmargin=2em]
\item 让错误信息转化为下一轮可利用的显式经验；
\item 对代码、查询、搜索类任务特别有效；
\item 可把历史失败模式沉淀为 procedural memory（程序性记忆）。
\end{itemize}

但注意，Reflexion 不是``让模型空想''。好的反思必须绑定\textbf{可验证反馈}，比如测试失败日志、HTTP 403、schema validation error、单元测试断言差异等。


\section{8.7 工具使用（Tool Use）设计}

Agent 能不能稳定工作，70\% 取决于工具层设计。


\subsection{8.7.1 工具描述的最佳结构}

推荐统一格式：


\begingroup
\scriptsize
\setlength{\tabcolsep}{3pt}
\begin{longtable}{|p{0.31\textwidth}|p{0.31\textwidth}|p{0.31\textwidth}|}
\hline
字段 & 说明 & 设计建议 \\ \hline
name & 短小唯一 & 动词开头，如 \texttt{search\_docs} \\ \hline
description & 告诉模型何时用、何时不用 & 写清楚适用边界 \\ \hline
parameters & 严格 schema & 限制类型、范围、枚举 \\ \hline
examples & 可选 & 给 1\textasciitilde{}2 个典型调用 \\ \hline
permissions & 可选 & 标注 read-only / write \\ \hline
\end{longtable}
\endgroup

差的描述：``查询信息。''


好的描述：``在公司知识库中搜索技术文档。适用于查找内部 API、部署流程、架构设计；不适用于互联网公开搜索。''


\subsection{8.7.2 模型如何选择工具}

模型本质上会做一个隐式分类问题：


\begin{quote}
当前用户意图更适合直接回答，还是调用哪一个工具？\\
\end{quote}

因此你要让工具集合具备：


\begin{itemize}[leftmargin=2em]
\item \textbf{低歧义}：避免同时有 \texttt{search\_web}、\texttt{web\_lookup}、\texttt{internet\_search} 三个几乎一样的工具；
\item \textbf{高可分性}：工具边界清晰；
\item \textbf{有限数量}：常见经验是 10\textasciitilde{}30 个核心工具比 100 个泛工具更稳定。
\end{itemize}

\subsection{8.7.3 并行工具调用}

并行（parallel tool calling）非常关键。比如用户问：


``比较北京、上海、深圳今天温度，并给出最高和最低城市。''


如果串行查三次天气，整体延迟可能是：


\begin{itemize}[leftmargin=2em]
\item 每次 API 400ms
\item 串行 = 1.2s + LLM 推理
\item 并行 = 0.4s + 聚合
\end{itemize}

并行调用架构：


\begin{codeblock}
            +--> get_weather(Beijing) --+
User Query -+--> get_weather(Shanghai) -+--> Aggregator --> Final Answer
            +--> get_weather(Shenzhen) -+
\end{codeblock}

适合并行的前提：


\begin{itemize}[leftmargin=2em]
\item 工具间无依赖；
\item 结果聚合规则明确；
\item 成本可控；
\item 下游能处理部分失败。
\end{itemize}

\subsection{8.7.4 工具结果处理与错误恢复}

工具执行不只是``返回字符串''。推荐统一返回结构：


\begin{codeblock}
{
    "ok": True,
    "data": {...},
    "error": None,
    "retryable": False
}
\end{codeblock}

错误恢复策略：


\begin{enumerate}[leftmargin=2em]
\item 参数校验失败：把 schema error 回传模型，让其自修正；
\item 临时网络失败：自动重试 2\textasciitilde{}3 次；
\item 权限不足：立即终止该工具并提示用户；
\item 非关键工具失败：降级回答；
\item 关键路径失败：触发 re-plan。
\end{enumerate}

\section{8.8 规划策略}

\subsection{8.8.1 Chain-of-Thought Planning}

Chain-of-Thought（思维链）规划适合短任务。模型先在内部展开若干步思考，再给出行动。它的优点是简单，缺点是对长任务不稳定，而且不适合完全暴露给用户。


\subsection{8.8.2 递归任务分解}

工程上更有价值的是递归分解：


\begin{codeblock}
Goal: 生成竞品分析
 ├─ 收集竞品名单
 ├─ 获取每个竞品最近版本信息
 │   ├─ 读取官网
 │   ├─ 搜新闻
 │   └─ 提取功能
 ├─ 汇总对比表
 └─ 输出建议
\end{codeblock}

递归分解的核心不是``拆得越细越好''，而是把每个子任务拆到：


\begin{itemize}[leftmargin=2em]
\item 输入明确；
\item 工具明确；
\item 完成标准明确；
\item 可独立验证。
\end{itemize}

\subsection{8.8.3 目标导向规划与回溯}

当任务存在障碍时，Agent 需要 backtracking（回溯）。例如：


\begin{enumerate}[leftmargin=2em]
\item 原计划通过 API 获取价格；
\item API 权限不足；
\item 回退到公开文档抓取；
\item 若仍失败，再输出``基于现有公开数据''的近似结果。
\end{enumerate}

这和搜索算法里的 DFS/BFS 不完全一样，但思想相通：\textbf{不是失败就结束，而是换路径继续逼近目标。}


\section{8.9 Agent Loop 中的状态管理}

没有状态管理的 Agent，通常活不过 5 轮。


最少需要管理三类状态：


\begingroup
\scriptsize
\setlength{\tabcolsep}{3pt}
\begin{longtable}{|p{0.31\textwidth}|p{0.31\textwidth}|p{0.31\textwidth}|}
\hline
状态 & 作用 & 示例 \\ \hline
Conversation State & 保存对话与工具历史 & messages, tool traces \\ \hline
Task State & 保存当前目标和计划 & 当前 step=3/5 \\ \hline
Environment State & 保存外部执行上下文 & 文件路径、session id、权限级别 \\ \hline
\end{longtable}
\endgroup

常见工程问题：


\begin{itemize}[leftmargin=2em]
\item 循环次数上限如何设？常见 6\textasciitilde{}20 步；
\item 何时终止？可由模型输出 \texttt{finish}，也可由外部策略判断；
\item 工具结果要不要全文写回上下文？通常只写摘要和关键字段；
\item 如何避免重复调用同一工具？给状态加去重缓存。
\end{itemize}

一个实用状态对象可能包含：


\begin{codeblock}
state = {
    "goal": "...",
    "history": [],
    "observations": [],
    "plan": [],
    "current_step": 0,
    "tool_cache": {},
    "failures": [],
    "finished": False
}
\end{codeblock}

\section{8.10 架构模式对比：什么时候用哪一种}

\begingroup
\scriptsize
\setlength{\tabcolsep}{3pt}
\begin{longtable}{|p{0.19\textwidth}|p{0.19\textwidth}|p{0.19\textwidth}|p{0.19\textwidth}|p{0.19\textwidth}|}
\hline
模式 & 适合场景 & 优势 & 劣势 & 面试回答关键词 \\ \hline
ReAct & 通用问答、检索、轻执行 & 简单、可解释、易实现 & 长任务效率一般 & trace、interleaving \\ \hline
Function Calling & 企业 API、结构化工具 & 稳定、可验证、强约束 & 不是完整规划框架 & schema、structured call \\ \hline
Plan-and-Execute & 复杂多步任务 & 全局性更强、便于审计 & 多一次规划成本 & planner/executor \\ \hline
Reflexion & 需要试错纠错 & 利于自改进 & 额外 token 成本 & self-evaluation、feedback \\ \hline
\end{longtable}
\endgroup

一个很好用的面试表达是：


\begin{quote}
小任务用 ReAct，中任务用 Function Calling + ReAct，大任务用 Plan-and-Execute，失败率高的任务叠加 Reflexion。\\
\end{quote}

\section{8.11 从零实现一个简单 ReAct Agent}

下面我们不用任何框架，只用 Python + OpenAI API 风格接口，实现一个最小可运行 ReAct Agent。为了让代码自包含，示例工具使用本地函数模拟搜索与计算。


\subsection{8.11.1 项目结构}

\begin{codeblock}
react_agent/
├─ agent.py
└─ requirements.txt
\end{codeblock}

\texttt{requirements.txt}


\begin{codeblock}
openai>=1.40.0
\end{codeblock}

\subsection{8.11.2 完整代码}

\begin{codeblock}
import json
import math
from typing import Any, Dict, List

from openai import OpenAI


client = OpenAI(api_key="YOUR_API_KEY")


def calculator(expression: str) -> Dict[str, Any]:
    allowed = {
        "abs": abs,
        "round": round,
        "sqrt": math.sqrt,
        "pow": pow,
    }
    try:
        value = eval(expression, {"__builtins__": {}}, allowed)
        return {"ok": True, "data": {"expression": expression, "result": value}, "error": None}
    except Exception as exc:
        return {"ok": False, "data": None, "error": str(exc)}


def fake_search(query: str) -> Dict[str, Any]:
    corpus = {
        "python release date": "Python 3.12 was released on 2023-10-02.",
        "openai founded": "OpenAI was founded in 2015.",
        "anthropic founded": "Anthropic was founded in 2021."
    }
    result = corpus.get(query.lower(), "No exact result found.")
    return {"ok": True, "data": {"query": query, "result": result}, "error": None}


TOOLS = {
    "calculator": {
        "description": "执行数学表达式计算，适用于加减乘除、平方根、幂运算。",
        "schema": {
            "type": "object",
            "properties": {
                "expression": {"type": "string", "description": "例如 (32 + 18) / 5"}
            },
            "required": ["expression"],
            "additionalProperties": False,
        },
        "handler": calculator,
    },
    "search": {
        "description": "搜索小型知识库，适用于查询已知事实。",
        "schema": {
            "type": "object",
            "properties": {
                "query": {"type": "string", "description": "搜索关键词"}
            },
            "required": ["query"],
            "additionalProperties": False,
        },
        "handler": fake_search,
    },
}


SYSTEM_PROMPT = """
你是一个 ReAct Agent。
你必须严格遵循以下格式之一：

1. 如果需要调用工具，输出 JSON：
{"type":"tool_call","name":"工具名","arguments":{...},"thought":"简短说明"}

2. 如果已经可以回答，输出 JSON：
{"type":"final","answer":"最终答案","thought":"简短说明"}

不要输出额外文本。
"""


def build_tool_spec() -> str:
    specs = []
    for name, meta in TOOLS.items():
        specs.append({
            "name": name,
            "description": meta["description"],
            "parameters": meta["schema"],
        })
    return json.dumps(specs, ensure_ascii=False, indent=2)


def call_model(messages: List[Dict[str, str]]) -> Dict[str, Any]:
    response = client.chat.completions.create(
        model="gpt-4.1-mini",
        temperature=0,
        messages=messages,
        response_format={"type": "json_object"},
    )
    content = response.choices[0].message.content
    return json.loads(content)


def run_agent(user_goal: str, max_steps: int = 8) -> str:
    messages: List[Dict[str, str]] = [
        {"role": "system", "content": SYSTEM_PROMPT},
        {
            "role": "user",
            "content": (
                f"可用工具如下：\n{build_tool_spec()}\n\n"
                f"用户目标：{user_goal}"
            ),
        },
    ]

    for step in range(1, max_steps + 1):
        decision = call_model(messages)
        print(f"[Step {step}] Decision => {decision}")

        if decision["type"] == "final":
            return decision["answer"]

        if decision["type"] != "tool_call":
            raise ValueError(f"Unknown decision type: {decision}")

        tool_name = decision["name"]
        arguments = decision["arguments"]

        if tool_name not in TOOLS:
            observation = {
                "ok": False,
                "error": f"Tool {tool_name} not found",
                "retryable": False,
            }
        else:
            try:
                observation = TOOLS[tool_name]["handler"](**arguments)
            except TypeError as exc:
                observation = {
                    "ok": False,
                    "error": f"Bad arguments: {exc}",
                    "retryable": True,
                }

        messages.append(
            {
                "role": "assistant",
                "content": json.dumps(decision, ensure_ascii=False),
            }
        )
        messages.append(
            {
                "role": "user",
                "content": "Observation: " + json.dumps(observation, ensure_ascii=False),
            }
        )

    return "超过最大执行步数，任务未完成。"


if __name__ == "__main__":
    answer = run_agent("OpenAI 成立于哪一年？再计算 2026 - 该年份。")
    print("Final Answer:", answer)
\end{codeblock}

\subsection{8.11.3 代码解读}

这个实现虽然小，但具备了 ReAct 的关键要素：


\begin{enumerate}[leftmargin=2em]
\item \textbf{显式循环}：\texttt{for step in range(...)}；
\item \textbf{工具注册表}：\texttt{TOOLS}；
\item \textbf{结构化决策}：模型只能输出 \texttt{tool\_call} 或 \texttt{final}；
\item \textbf{Observation 回写}：工具结果重新注入上下文；
\item \textbf{最大步数保护}：避免死循环。
\end{enumerate}

如果要扩展到生产环境，下一步通常是：


\begin{itemize}[leftmargin=2em]
\item 用真正的 function calling 替代 JSON 文本协议；
\item 加入参数校验；
\item 接入日志与 trace id；
\item 对工具结果做摘要压缩；
\item 增加重试、缓存与权限控制；
\item 把状态持久化到 Redis 或数据库。
\end{itemize}

\section{8.12 面试中如何讲 Agent 架构}

如果面试官问：``请你设计一个企业知识库问答 Agent 架构''，你可以按以下顺序回答：


\begin{enumerate}[leftmargin=2em]
\item 明确目标：问答、检索、引用、权限控制；
\item 选模式：用 Function Calling + ReAct；
\item 状态：保留会话历史、工具 trace、检索缓存；
\item 工具：搜索、文档读取、权限校验；
\item 错误恢复：检索无结果时扩大召回，工具失败时重试；
\item 可观测性：记录每步 thought summary、tool name、latency、token usage；
\item 安全：只读工具、敏感文档 ACL。
\end{enumerate}

这样的回答会让人感觉你不是``会调模型''，而是在设计一套可上线系统。


\section{8.13 生产中如何组合这些模式}

实际项目里很少只用一种模式。更常见的做法是把几种模式叠加成分层架构：


\begin{enumerate}[leftmargin=2em]
\item 最外层用 \textbf{Plan-and-Execute} 做长任务编排；
\item 每个执行步骤内部用 \textbf{ReAct} 做短循环；
\item 工具层统一走 \textbf{Function Calling}；
\item 遇到失败率高的步骤，再插入 \textbf{Reflexion} 做自我修正。
\end{enumerate}

一个真实感很强的例子是``代码仓库分析 Agent''：


\begin{itemize}[leftmargin=2em]
\item 用户目标：分析某仓库中最近 30 天的性能回归并给出修复建议；
\item Planner 先拆成``找 commit、定位回归、读取性能日志、生成建议''四步；
\item `\texttt{定位回归''这一步内部，Executor 用 ReAct 反复调用 }git\_log\texttt{、}read\_file\texttt{、}search\_trace`；
\item 工具调用全部使用结构化 schema；
\item 如果生成的根因解释与测试日志矛盾，则启动 Reflexion，再读一次关键信息。
\end{itemize}

这种组合方式的好处，是把`\texttt{全局规划''和}`局部试错''分开。全局层追求正确方向，局部层追求灵活执行。如果把两者混在一起，模型很容易在细枝末节上花掉全部上下文预算。


再看一个适用于面试的选型表：


\begingroup
\scriptsize
\setlength{\tabcolsep}{3pt}
\begin{longtable}{|p{0.31\textwidth}|p{0.31\textwidth}|p{0.31\textwidth}|}
\hline
任务类型 & 推荐组合 & 原因 \\ \hline
查询天气、库存、订单 & Function Calling + ReAct & 步骤短，重点是工具稳定性 \\ \hline
文档总结、资料调研 & ReAct + Summary Memory & 检索后快速整合即可 \\ \hline
数据分析、代码修复 & ReAct + Reflexion & 错误反馈可验证，适合自修正 \\ \hline
长流程自动化 & Plan-and-Execute + Function Calling & 先规划、后分步执行更稳 \\ \hline
高风险写操作 & Plan-and-Execute + Human Approval & 需要检查点与审批点 \\ \hline
\end{longtable}
\endgroup

一个成熟工程师与初学者的差别，往往不在于会不会背这些名字，而在于能否说清楚``为什么这个场景要这样组合''。


\section{8.14 Agent 架构的关键监控指标}

如果系统要上线，就不能只看``回答像不像人''。至少要埋以下指标：


\begingroup
\scriptsize
\setlength{\tabcolsep}{3pt}
\begin{longtable}{|p{0.31\textwidth}|p{0.31\textwidth}|p{0.31\textwidth}|}
\hline
指标 & 含义 & 典型阈值 \\ \hline
loop\_steps & 单次任务循环步数 & 6\textasciitilde{}12 步为常见范围 \\ \hline
tool\_success\_rate & 工具调用成功率 & 核心工具应 > 95\% \\ \hline
repeated\_action\_rate & 重复调用同一工具的比例 & 持续升高通常表示陷入局部循环 \\ \hline
final\_answer\_rate & 能正常收敛到最终答案的比例 & 生产目标通常 > 90\% \\ \hline
avg\_tokens\_per\_task & 单任务 token 消耗 & 用于成本治理 \\ \hline
replanning\_rate & 需要重规划的比例 & 过高说明 planner 质量不足 \\ \hline
\end{longtable}
\endgroup

其中最容易被忽略的是 \texttt{repeated\_action\_rate}。很多 Agent 表面上``没有报错''，但它会连续三次调用同一个搜索工具，只是换着措辞问。这类问题如果不靠日志统计，很难单凭肉眼发现。


另一个关键点是 \textbf{工具观察结果的可压缩性}。如果某个工具每次返回 20KB 原始 JSON，而模型真正需要的只有两个字段，那么系统应该在 Observe 阶段就做 reducer，把结果压成：


\begin{codeblock}
订单状态=已发货；预计送达=2026-06-14；物流单号=SF123456789
\end{codeblock}

这样既减少 token，又能提升后续推理质量。


\section{8.15 三个常见反模式}

\subsection{反模式一：把 Thought 当成对用户的产品功能}

有些系统把模型的内部思考原样展示给用户，结果既冗长又不稳定。正确做法是：\textbf{内部 trace 用于调试，面向用户的解释单独生成}。


\subsection{反模式二：工具层太薄，只返回``成功/失败''}

如果工具错误只返回 \texttt{failed}，模型根本不知道如何修。更好的返回应该告诉它：


\begin{itemize}[leftmargin=2em]
\item 哪个字段错了；
\item 期望类型是什么；
\item 是否可重试；
\item 是否建议改用别的工具。
\end{itemize}

\subsection{反模式三：没有显式终止条件}

很多新人写 Agent Loop 时，只让模型自己判断何时结束，没有外部硬限制。结果一旦模型进入循环，就会一直消耗 token。实际工程中至少要同时有三道闸：


\begin{enumerate}[leftmargin=2em]
\item 最大循环步数；
\item 最大工具调用次数；
\item 最大 token / 成本预算。
\end{enumerate}

你在面试里如果能主动提这三点，面试官通常会认为你做过真实系统，而不是只跑过 notebook。


\section{8.16 一个实战设计题的回答框架}

假设面试官让你现场设计``一个能够读取企业文档、查询工单系统、自动生成处理建议的客服 Agent''。你可以直接按下面框架展开：


第一步，说明为什么它不是简单问答，而是 Agent：因为它要根据用户问题自主决定查知识库、查工单还是结束回答，还要在多步调用后汇总证据。


第二步，说明架构模式：外层使用 Function Calling + ReAct，复杂问题再加一个轻量 planner。理由是客服场景大部分问题在 3\textasciitilde{}6 步内可收敛，不一定需要重型 Plan-and-Execute，但又必须具备显式工具调用能力。


第三步，说明状态：


\begin{itemize}[leftmargin=2em]
\item 会话状态保存最近问答；
\item 任务状态保存当前是否已查工单、是否已查知识库；
\item 环境状态保存用户身份、权限、工单号。
\end{itemize}

第四步，说明终止条件：当模型拿到足够证据并且置信度达到阈值，比如 0.8，就输出最终答案；否则转人工。


第五步，说明监控：统计平均循环数、工具成功率、转人工率、重复调用率、引用证据覆盖率。


这种回答方式的价值在于，它把`\texttt{模式名词''落到了}`如何上线''。


\section{本章要点}

\begin{itemize}[leftmargin=2em]
\item Agent 与单次 LLM 调用的根本区别在于自主性、工具使用、规划和记忆。
\item 绝大多数 Agent 都可抽象成 Perceive -> Reason -> Act -> Observe 的循环。
\item ReAct 适合通用场景，优点是简单、透明、易调试。
\item Function Calling 把工具调用接口结构化，是现代 Agent 工程的基础设施。
\item Plan-and-Execute 适合复杂长任务，Reflexion 适合高失败率、可验证反馈场景。
\item 工具设计要强调 schema、边界、并行能力、错误恢复和统一返回格式。
\item 没有状态管理与终止条件的 Agent，很容易失控、重复调用或无限循环。
\end{itemize}

\section{延伸阅读}

\begin{enumerate}[leftmargin=2em]
\item ReAct: Synergizing Reasoning and Acting in Language Models
\item Toolformer: Language Models Can Teach Themselves to Use Tools
\item OpenAI Function Calling / Responses API 官方文档
\item Anthropic Tool Use 与 Model Context Protocol 官方文档
\item LangGraph 关于 stateful agent loops 的设计文档
\end{enumerate}
